\section{Introduction}\label{b:introduction}

The spacetime Penrose inequality compares the invariant ADM mass
$m_{\mathrm{ADM}}=\sqrt{E^2-|P|^2}$ with the area needed to enclose a trapped
region, where $E$ and $P$ are the energy and linear momentum of
asymptotically flat initial data.  The deformations below control energy
in a chosen end; for general weak data, the invariant-mass conclusion
uses a separate passage to rest ends, where $P=0$.
Reducing the problem to the Riemannian Penrose inequality already
exposes the main difficulty.  A deformation must create nonnegative scalar
curvature while controlling both quantities in the comparison: energy
at infinity and the area of every enclosing cut.  If the deformation is
performed on an exterior with a genuine inner boundary, its boundary
flux is a third quantity that must be controlled.

This paper constructs such boundary deformations in three and four
spatial dimensions.  It also develops a direct four-dimensional maximal
construction that retains a noncompact decaying second fundamental form.
The common output is a smooth complete Riemannian exterior with
nonnegative scalar curvature, strictly negative inner mean curvature,
a metric lower bound protecting all enclosing areas, and an arbitrarily
small upper error in energy.  A minimal enclosure then supplies the
boundary to which the Riemannian Penrose inequality applies.

To explain the area comparison, let $(M,g,K)$ be an initial data exterior
with compact boundary and one asymptotically Euclidean end.  The tensor
$K$ is the second fundamental form.  Orient an inner boundary toward
the end, let $H$ be its mean curvature, and write
$\theta_+=H+\tr_{\mathrm{tan}}K$ for its future expansion.
A full enclosing cut is the entire compact boundary of a connected
closed outer domain that contains the distant end and has interior in
the open exterior.  Every boundary component counts, including any
part coincident with the original boundary.  Denote the infimum of their
$g$-areas by $A_*(g)$.  The infimum need not be attained.
The precise asymptotic and constraint hypotheses for each construction
are stated at its beginning.

Both the area and the horizon predicate matter. In spatial dimension
three, Ben-Dov's example
has an outermost future apparent horizon $S$ with
$|S|>16\pi m_{\mathrm{ADM}}^2$~\cite{BenDov2004}.
Carrasco--Mars construct outer area-minimizing generalized apparent
horizons, defined by $H=|\tr_{\mathrm{tan}}K|$, with the same strict area
inequality~\cite{CarrascoMars2010}.
These concern distinct formulations.  We use full enclosing area and
specify the expansion condition in each theorem.

The area comparison uses an elementary point that remains valid for all
these cuts.  In spatial dimension $d$, if $h\ge c g$ as quadratic forms
for some $c>0$, then the restriction to each tangent $(d-1)$-plane gives
\[
 |\Gamma|_h\ge c^{(d-1)/2}|\Gamma|_g
 \quad\hbox{and hence}\quad
 A_*(h)\ge c^{(d-1)/2}A_*(g).
\]
No selection of a minimizing component and no convergence of minimizers
is needed for this inequality.  We use it after constructing a metric
of the form
\[
 \widehat g=e^{4t/(d-2)}\bigl(g+l(t)^2\dd f^2\bigr),
 \qquad t\ge-\varepsilon.
\]
Here $f$ is the graph height, $t$ is the conformal unknown, and $l>0$
is the warping factor.
The graph term is nonnegative on every tangent plane.  Thus the lower
bound for $t$ protects the full enclosing area before the graph is
known to have bounded slope.

\subsection{The constructions and their relationship}

The first two constructions start with strict dominant-energy data,
a strictly future-trapped boundary, and a prepared end with compactly
supported $K$. The third keeps a maximal tensor tail. The following
map separates these geometric outputs from their original-data
applications; the cited statements give the full hypotheses.
Here $\omega_3=|S^3|=2\pi^2$.

\begin{center}
\small
\renewcommand{\arraystretch}{1.2}
\begin{tabular}{@{}p{.25\linewidth}p{.30\linewidth}p{.37\linewidth}@{}}
\toprule
Construction & Geometric output & Original-data consequence \\
\midrule
Prepared dimension three; compact $K$
& Theorem~\ref{n3:prepared:output}
& $m_{\mathrm{ADM}}\ge\sqrt{A_*/(16\pi)}$ for weak decay $q>1/2$;
Theorem~\ref{n3:intro:exterior-inequality} \\[3pt]
Prepared dimension four; compact $K$
& Theorem~\ref{four:thm:boundary-deformation}
& $m_{\mathrm{ADM}}\ge\frac12(A_*/\omega_3)^{2/3}$ for weak decay $q>1$;
Corollary~\ref{b:four-weak} \\[3pt]
Direct maximal dimension four; decaying $K$
& Proposition~\ref{four:maximal:prop:geometric-output}
& $E\ge\frac12(|S|_g/\omega_3)^{2/3}$ in the maximal-vacuum application;
Theorem~\ref{four:maximal:thm:main} \\
\bottomrule
\end{tabular}
\end{center}

In the first two applications, weak decay means
$g-\delta=O_2(r^{-q})$ and $K=O_1(r^{-1-q})$.
The exteriors satisfy the dominant energy condition, integrable
constraint densities, finite timelike ADM charges $E>|P|$, and weak
future trapping. Their separate rest preparations control all enclosing
cuts. In dimension three, the prepared energies tend to
$m_{\mathrm{ADM}}$ and the full enclosing-area infima converge. In
dimension four, the distant-end replacement supplies a one-sided lower
area bound; strictification is removed at each fixed repaired end before
the replacement radius tends to infinity. These inputs are
\cite[Proposition~7.13 and Proposition~8.8/Lemma~8.9, respectively]{NumericalCompanion}.

The direct maximal theorem has its own approximation and retains the
tensor tail. Its stated application has $q>1$, ball-exterior topology,
a connected spherical marginally outer trapped boundary ($\theta_+=0$),
a torus action, $P=0$, $E>0$,
outermostness, and outer area minimization against all enclosures.
Here $P=0$ identifies energy with invariant mass, and outer area
minimization identifies $A_*$ with $|S|_g$. The strict geometric output
in the middle column has its own hypotheses; the theorem records this
particular original-data application. The broader numerical comparison
in Corollary~\ref{b:four-weak} uses the distinct end-replacement route.

The difference at infinity is concrete.  With compactly supported $K$,
the distant trace equation is a homogeneous graph equation with a
positive height term.  When $\tr_gK=0$, the retained tensor contributes a
term quadratic in $\nabla f$, which becomes a drift along the solution.
The equation remains homogeneous and still forces exponential graph
decay at fixed deformation parameter.  The scalar equation behaves
differently: its zero-graph limit retains $|K|^2/2$.
Section~\ref{b:transition} derives this term and the tail weights needed
to control it.  That calculation identifies both the shared argument
and the additional preparation, threshold geometry and estimates required
for the maximal route.

\subsection{Why the inner boundary matters}

The scalar-curvature identity separates the original constraint densities,
a nonnegative quadratic expression, and the divergence of a vector field
$V$.  Prescribing this divergence produces the required curvature sign.
At infinity, the integral of $V$ measures the ADM energy change.  A
boundary condition for $V$ also prescribes the mean curvature of the
deformed inner boundary.  The divergence theorem therefore links the
two ends of the construction.

The useful boundaries are threshold trapped-region frontiers.  For
general $K$, two signs of the tensor produce black and white faces,
whose names distinguish their expansion equations and the corresponding
Dirichlet heights of the graph.  The collars outside these faces supply
geometric barriers for the graph before uniform ellipticity is available.
They first confine its scaled height to the required interval and then
force a signed normal slope.  That slope makes the possible negative
inner flux exponentially small relative to the positive bulk terms.
The all-black maximal construction has its own threshold argument,
using maximality in place of the two-sided trace budget.

This order is essential.  Uniform ellipticity is an outcome of the
geometric height, floor and slope estimates; it is not assumed to prove
them.  Once the coefficients lie in a bounded elliptic range, local
scalar estimates give the boundary regularity needed for a compact
continuation map.  The four-dimensional proof retains the critical
interpolation step and the interface measures produced by reflection.
Each application verifies its own boundary conditions and completes
its own degree and bootstrap argument.

There are two kinds of constants throughout.  Polynomial bounds in the
large deformation parameter $N$ are used in the final flux comparison.
Classical regularity constants may depend arbitrarily on fixed $N$.
One first exhausts the outer radius with $N$ fixed, then lets $N$ tend
to infinity using only the polynomial estimates.  This distinction
prevents a local compactness argument from being mistaken for a
uniform estimate in the final mass limit.

\subsection{Relation to earlier methods}

Penrose's proposed inequality arose from the relation between gravitational
collapse, horizons and total energy~\cite{Penrose1973}.  In the
time-symmetric setting the second fundamental form vanishes and the
dominant-energy condition becomes nonnegative scalar curvature.
Geroch's smooth-flow idea and the conditional argument of Jang--Wald
preceded the weak-flow theory~\cite{Geroch1973,JangWald1977}; see also
the historical discussion in~\cite[p.~359]{HuiskenIlmanen2001}.
Huisken--Ilmanen's weak inverse-mean-curvature flow proved the
three-dimensional Riemannian inequality for each connected component
of the horizon~\cite[Main Theorem, p.~356]{HuiskenIlmanen2001};
Bray's conformal flow gives the bound for the total area of a disconnected
horizon~\cite[Theorems~1 and~19]{Bray2001}.  Bray--Lee extended the numerical
conformal-flow theorem to dimensions below eight
\cite[Theorem~1.4]{BrayLee2009}.  The final comparison
here uses the numerical clauses of these theorems after constructing
their complete Riemannian inputs.

The graph strategy comes from Jang's equation and the Schoen--Yau
positive-energy argument~\cite{Jang1978,SchoenYau1981}. The earlier
Schoen--Yau proof establishes the positive mass theorem in the Riemannian
and maximal-data cases by minimal surfaces~\cite{SchoenYau1979}.  A classical Jang graph can develop
a cylindrical end at an apparent horizon; Metzger analyses this blow-up
and its exponential cylindrical convergence near a strictly stable
component~\cite{Metzger2010}.
Conformal normalization of such a graph can lose the horizon-area
comparison.  Bray and Khuri developed the warped-graph approach to
address this obstruction~\cite{BrayKhuri2010,BrayKhuri2011}; their 2010
work also proves a sharp spherical case under its stated hypotheses.
The general warped scalar identity already retains a trace discrepancy
before the generalized Jang equation is imposed
\cite[Identity~9, p.~580]{BrayKhuri2011}.

Han--Khuri prove existence and quantitative boundary blow-up for
prescribed warping factors~\cite{HanKhuri2013}.  Their later coupling
to conformal flow gives a reduction contingent on solving the coupled
system with the required boundary behavior~\cite{HanKhuri2018}.
The direct conclusion of that argument concerns ADM energy, as its
Remark~2 explains.  Here the coupled equations prescribe a nonzero
divergence source, a trace depending increasingly on graph height, and
a penalty near the conformal lower bound.  The boundary analysis proves
global existence for these equations and supplies the area and energy
comparisons needed for the subsequent limits.

This distinction concerns the equations, not only their interpretation.
Jaracz proves radial nonexistence for a specified generalized-Jang and
zero-divergence system~\cite{Jaracz2023}; that result does not apply to
the nonzero-source system studied here. His spacetime Poisson method
also provides a predecessor for conformal strictification
\cite{Jaracz2025StrictDEC}. The maximal construction below solves its
own problem with nonzero inner Neumann data and proves the energy and
full enclosing-area comparisons needed to remove the strict
approximation.

Trapped-region boundaries and their stability enter through the work of
Andersson--Metzger and Andersson--Eichmair--Metzger
\cite{AnderssonMetzger2009,AnderssonEichmairMetzger2011}.
Eichmair's Perron and almost-minimizing constructions provide the
original higher-dimensional method~\cite{Eichmair2010}; the precise
total-region theorem used below is the one recorded by
Andersson--Eichmair--Metzger.  For the stability operator of marginally
outer trapped surfaces (MOTS) and its principal-eigenfunction
characterization of stability we use
Andersson--Mars--Simon~\cite{AnderssonMarsSimon2008}.  The local arguments
below explain the additional threshold continuity, support transport
and collar choices used by the deformation.  The analytic part combines
scalar elliptic estimates with the Leray--Schauder degree principle
\cite{LeraySchauder1934}; the particular boundary homotopies and their
admissible open sets are constructed explicitly.

Recent spacetime results also distinguish the relevant scopes.
Khuri--Kunduri prove the sharp inequality for cohomogeneity-one
$\mathrm{SU}(n+1)$-invariant data in spatial dimensions $2(n+1)\ge4$,
with their outermost future-or-past horizon condition
\cite[Theorem~1.1]{KhuriKunduri2025}.  The direct maximal theorem here
has different symmetry and outer-area-minimization hypotheses.
Allen--Bryden--Kazaras--Khuri prove a three-dimensional invariant-mass
bound by a universal suboptimal constant times the square root of
minimum enclosing area, with additional trace decay and a specified
horizon selection~\cite[Theorem~1.1 and Remark~1.2]{AllenBrydenKazarasKhuri2025}.

Two companion manuscripts supply specified inputs: scalar algebra and
weak-end preparations from~\cite{NumericalCompanion}, and full-perimeter
enclosure geometry from~\cite{RiemannianCompanion}. Their exact hypotheses
and normalization are checked at each use. The numerical comparisons
at the end of the boundary constructions use Bray and Bray--Lee directly.

\subsection{Reading the proof}

The three routes retain their own geometry, equations, global bounds and
end comparison. Shared arguments are proved at their first use in the
three-dimensional route: threshold continuity, stable collars and exact
support transport in Section~\ref{n3:domains:section}; the passage from
lower tests to sublevel supports in Lemma~\ref{n3:apriori:sublevel-support};
and scalar measure-source regularity and degree on varying open sets in
Lemmas~\ref{b:lem:natural-growth} and~\ref{b:lem:moving-degree}.
Each later use verifies its tensor, coefficients and boundary conditions.

For dimension three, Sections~\ref{n3:deformation:section}--\ref{n3:limit:section}
proceed from the coupled equations and conformal floor through geometric
bounds, boundary existence, and the complete-end comparison. Dimension
four follows the same order in
Sections~\ref{four:sec:geometry}--\ref{four:sec:existence}, with its own
compact total-region construction and the critical interpolation in
Theorem~\ref{four:thm:axial-regularity}. Corollary~\ref{b:four-weak}
then gives its separate weak-decay transfer.

Section~\ref{b:transition} explains the maximal tensor tail before
Sections~\ref{four:maximal:sec:preparation}--\ref{four:maximal:sec:comparison}
construct the direct route. Its pointwise scalar and boundary identities
come from Section~\ref{four:sec:system}, its axial estimate from
Theorem~\ref{four:thm:axial-regularity}, and its enclosure input from
Proposition~\ref{four:input:enclosure}. The maximal part supplies its own
strict approximation, noncompact drift bounds, scalar solve, reflected
boundary checks and final flux estimate.
